\documentclass{article}

\usepackage{hyperref,natbib,graphicx,caption,amsmath,wasysym}

\title{Carbonate ${}^{206}$Pb/${}^{238}$U problems and potential ${}^{207}$Pb/${}^{235}$U fixes}

\author{P.~Vermeesch$^1$, N.~McLean$^2$, A.~Vaks$^3$, T.~Golan$^3$, S.F.M~Breitenbach$^4$ and
  R.~Parrish$^5$\\
  $^1$ University College London, London WC1E 6BT, United Kingdom\\
  $^2$ University of Kansas, Lawrence, KS 66045, United States\\
  $^3$ Geological Survey of Israel, 9692100 Jerusalem, Israel\\
  $^4$ Northumbria University, Newcastle upon Tyne NE1 8ST, United Kingdom\\
  $^6$ University of Portsmouth, Portsmouth PO1 3QL, United Kingdom
}

\begin{document}

\maketitle

\begin{abstract}
  Carbonate U--Pb dating has become a key tool for Quaternary
  palaeoclimatology and palaeoanthropology beyond the age limit of
  Th--U disequilibrium dating. U--Pb geochronology is based on the
  paired radioactive decay of ${}^{238}$U to ${}^{206}$Pb and of
  ${}^{235}$U to ${}^{207}$Pb.  Current carbonate U--Pb data
  processing algorithms rely mostly on the ${}^{206}$Pb/${}^{238}$U
  clock and attach little weight to the ${}^{207}$Pb/${}^{235}$U
  data. A key weakness of this approach is the need to correct the
  ${}^{206}$Pb/${}^{238}$U data for initial ${}^{234}$U/${}^{238}$U
  disequilibrium, which may cause an excess or (occasionally) a
  deficit in radiogenic ${}^{206}$Pb compared to secular
  equilibrium. Uncorrected initial disequilibrium may bias
  ${}^{206}$Pb/${}^{238}$U dates by up to 4\,Myr. We introduce a new
  disequilibrium correction algorithm, using matrix exponentials. This
  algorithm can be used to undo the effects of U-series disequilibrium
  using either an assumed initial composition, or a measured set of
  modern ${}^{234}$U/${}^{238}$U (and optionally
  ${}^{230}$Th/${}^{238}$U) activity ratios. Using a deterministic
  Bayesian inversion algorithm, we show that disequilibrium
  corrections work well for relatively young samples but become
  unreliable beyond 1.5\,Ma and impossible beyond 2\,Ma.  Using
  theoretical models and real world examples from Siberia, South
  Africa and Israel, we show that the uncertainty of the
  disequilibrium correction of such old samples exceeds the correction
  itself. Previous `Monte Carlo' error propagation methods
  underestimate these uncertainties by up to an order of magnitude.
  For carbonates older than 2\,Ma that likely experienced significant
  initial ${}^{234}$U/${}^{238}$U disequilibrium, we recommend using
  the ${}^{207}$Pb/${}^{235}$U isochron method instead of
  ${}^{206}$Pb/${}^{238}$U geochronology.  ${}^{207}$Pb/${}^{235}$U
  isochrons require only a small and simple correction for initial
  ${}^{231}$Pa depletion. This makes ${}^{207}$Pb/${}^{235}$U dating
  more accurate than ${}^{206}$Pb/${}^{238}$U geochronology. However,
  the ${}^{207}$Pb/${}^{235}$U method is no panacea. Its precision is
  limited by the lower abundance of ${}^{207}$Pb compared to
  ${}^{206}$Pb. In some samples, this loss of precision results in a
  failure to outperform the Bayesian credible intervals of the
  disequilibrium-corrected ${}^{206}$Pb/${}^{238}$U dates. Such
  samples remain undateable, unless prior information is available to
  constrain the initial ${}^{234}$U/${}^{238}$U activity ratios.
\end{abstract}

%\copyrightstatement{TEXT} %% This section is optional and can be used for copyright transfers.

\section{Introduction}\label{sec:intro}

Carbonate rocks are only a minor component of the continental
crust. However, their scientific importance far outweighs their
volumetric abundance. Biogenic carbonates and speleothems document the
history of life, and of Earth's climate and environment. To generate
detailed time series of past changes from these carbonate archives, an
accurate and precise chronological framework is essential. This
framework is anchored in absolute time using radiometric dating
methods, with two techniques commonly employed for this
purpose. ${}^{230}$Th/U dating is the method of choice for young
samples whose ${}^{230}$Th, ${}^{234}$U and ${}^{238}$U activities are
out of secular equilibrium \citep{kaufman1965,
  ludwig2003b}. ${}^{206}$Pb/${}^{238}$U dating is the default method
for older rocks, in which the secular equilibrium between
${}^{230}$Th, ${}^{234}$U and ${}^{238}$U has been restored
\citep{smith1989, roberts2020}.\medskip

Ironically, the absence of detectable ${}^{234}$U/${}^{238}$U
disequilibrium compromises the accuracy of the
${}^{206}$Pb/${}^{238}$U method. Any initial excess or deficit of
${}^{234}$U and ${}^{230}$Th affects the ${}^{206}$Pb/${}^{238}$U
ratio and, hence, the age estimate derived therefrom. In clean,
detritus-free carbonates, it is often safe to assume the absence of
initial ${}^{230}$Th. This assumption is not valid for ${}^{234}$U,
which can be enriched (or occasionally depleted) relative to
${}^{238}$U by physiochemical processes such as (1) $\alpha$-recoil
ejection and preferential leaching of ${}^{234}$U from
$\alpha$-damaged mineral sites, and (2) chemical fractionation between
preferentially oxidised ${}^{234}$U$^{6+}$ and ${}^{238}$U$^{4+}$
\citep{fleischer1982,porcelli2003}.\medskip

Corrections for initial ${}^{234}$U/${}^{238}$U disequilibrium can be
done by either assuming a specific initial ${}^{234}$U/${}^{238}$U
activity ratio, or by inferring the initial ratio from any measured
residual ${}^{234}$U/${}^{238}$U-disequilibrium \citep{richards1998,
  woodhead2006, wendt1985, engel2019}.  In Section~\ref{sec:diseq} of
this paper, we review the second approach using matrix
exponentials. We show that initial ${}^{234}$U/${}^{238}$U
disequilibrium can bias ${}^{206}$Pb/${}^{238}$U dates by up to
4\,Myr.\medskip

Current disequilibrium correction algorithms use a `Monte Carlo'
approach to propagate the errors. In Section~\ref{sec:MC} we will show
that this approach can underestimate the analytical uncertainties of
${}^{206}$Pb/${}^{238}$U dates by an order of magnitude when samples
are within a few permil of secular equilibrium (which typically
happens before c.\,2\,Ma). This observation undermines the results of
several published studies \citep{vermeesch2025}.\medskip

In Section~\ref{sec:Bayes} we introduce a deterministic Bayesian
approach to estimate the uncertainties of disequilibrium-corrected
${}^{206}$Pb/${}^{238}$U dates. We use this alternative algorithm to
show that beyond c.\,2\,Ma, disequilibrium-corrected
${}^{206}$Pb/${}^{238}$U dates are impractically imprecise, unless
highly enriched initial ${}^{234}$U/${}^{238}$U activity ratios can be
ruled out \emph{a priori}. The large uncertainty associated with the
${}^{206}$Pb/${}^{238}$U method degrades its ability to constrain
reliable chronologies for carbonates whose initial disequibrium has
expired. However, a more accurate approach is available. Following the
example of \citet{richards1998}, \citet{neymark2008}, \citet{vaks2020}
and others, Section~\ref{sec:207Pb235U} makes a case for the
little-used ${}^{207}$Pb/${}^{235}$U clock as a replacement for the
${}^{206}$Pb/${}^{238}$U method.\medskip

${}^{207}$Pb/${}^{235}$U isochrons are, essentially, immune to the
effects of initial disequilibrium (apart from a minor correction for
${}^{231}$Pa, which becomes smaller with increasing age).  In
Section~\ref{sec:accuracy-precision} we present examples from Siberia
and Israel to show that the ${}^{207}$Pb/${}^{235}$U method is more
accurate than the ${}^{206}$Pb/${}^{238}$U method, whilst being less
precise for young samples. Both the Bayesian uncertainty estimation
method and ${}^{207}$Pb/${}^{235}$U isochrons have been implemented in
the \texttt{IsoplotR} toolbox for radiometric geochronology
(Section~\ref{sec:IsoplotR}).\medskip

This paper will use the following symbols and notations:
\begin{itemize}
\item{$n_{38}$, $n_{34}$, $n_{30}$, $n_{26}$ and $n_{06}$}: the number
  of atoms of ${}^{238}$U, ${}^{234}$U, ${}^{230}$Th, ${}^{226}$Ra and
  ${}^{206}$Pb, respectively;
\item{$\lambda_{38}$, $\lambda_{35}$, $\lambda_{34}$, $\lambda_{32}$, $\lambda_{31}$ 
  $\lambda_{30}$ and $\lambda_{26}$}: the radioactive decay constants
  of ${}^{238}$U, ${}^{235}$U, ${}^{234}$U, ${}^{232}$Th, ${}^{231}$Pa, ${}^{230}$Th
  and ${}^{226}$Ra, respectively;
\item{$[4/8]_t$}: the present day ${}^{234}$U/${}^{238}$U activity
  ratio (i.e. $[\lambda_{34}\,n_{34}]/[\lambda_{38}\,n_{38}]$);
\item{$[4/8]_i$, $[0/8]_i$, $[1/5]_i$}: the initial
  ${}^{234}$U/${}^{238}$U, ${}^{230}$Th/${}^{238}$U and
  ${}^{231}$Pa/${}^{235}$U activity ratios, respectively;
\item{$[4/8]_m$ and $s[4/8]_m$}: the measured
  ${}^{234}$U/${}^{238}$U activity ratio and its standard error,
  respectively.
\item{$123.456 \pm 0.012$ represents a(n approximate) 95\% confidence
  interval; whereas 123.456(78) is a more succinct notation that
  indicates a value of 123.456 with a standard error of 0.078.}
\end{itemize}

\section{Disequilibrium corrections in a nutshell}\label{sec:diseq}

Even though the U--Pb decay systems consist of numerous steps (14 for
the $^{238}$U$\rightarrow^{206}$Pb chain and 11 for the
$^{235}$U$\rightarrow^{207}$Pb chain), conventional U--Pb
geochronology ignores this complexity and the method is mathematically
treated as a set of simple parent-daughter pairs.  This simplification
is justified once a state of secular equilibrium is established
between all the intermediate daughter products in the decay chains.
Such secular equilibrium is practically reached after 1 to 2 million
years.  As mentioned in Section~\ref{sec:intro}, any disequilibrium
that might exist prior to this secular equilibrium can be used as a
chronometer in its own right.\medskip

Initial disequilibrium of the U-decay series affects the accuracy of
the U--Pb method. For example, ignoring any initial excess
\textsuperscript{234}U results in an overestimated
\textsuperscript{206}Pb/\textsuperscript{238}U age, and ignoring any
initial \textsuperscript{231}Pa deficit results in an underestimated
\textsuperscript{207}Pb/\textsuperscript{235}U age. Therefore, initial
disequilibrium is one mechanism to produce discordant U--Pb results.
Extreme ${}^{234}$U-enrichments have been observed in places such as
South Africa \citep[\ensuremath{[4/8]_i<12};][]{kronfeld1994}, Siberia
\citep[\ensuremath{[4/8]_i<6};][]{vaks2020} and Japan
\citep[\ensuremath{[4/8]_i<12};][]{kuribayashi2025}. Using the
$[4/8]_i=12$ value as an upper bound, the maximum effect of initial
${}^{234}$U/${}^{238}$U-disequilibrium can be approximated as follows:
\begin{equation}
  \Delta(t) \approx \frac{1}{\lambda_{38}}\,
  \ln\left[
  1+\
  \frac{{}^{206}\mbox{Pb}}{{}^{238}\mbox{U}} +
  (12-1)\,\frac{\lambda_{38}}{\lambda_{34}} +
  (0-1)\,\frac{\lambda_{38}}{\lambda_{30}}
  \right] -
  \frac{1}{\lambda_{38}}\,\ln\left[1+\frac{{}^{206}\mbox{Pb}}{{}^{238}\mbox{U}}\right] \approx
  \frac{11}{\lambda_{34}} - \frac{1}{\lambda_{30}} = 3.8\,\mbox{Myr}
    \label{eq:worst-case}
\end{equation}

\noindent where $\lambda_{38}=0.155125(83),\mbox{Gyr}^{-1}$
\citep{jaffey1971}, $\lambda_{30}=9.1705(16)\,\mbox{Myr}^{-1}$ and
$\lambda_{34}=2.82206(80)\,\mbox{Myr}^{-1}$ \citep{cheng2013}. For old
carbonates ($>100$\,Ma, say), a 4\,Myr bias may be
inconsequential. However, for young carbonates, the relative effect of
initial disequilibrium can result in order-of-magnitude levels of
bias. A disequilibrium correction is needed to remove this
bias.\medskip

If the intermediate daughter is sufficiently long lived and the sample
is sufficiently young ($t < 5/\lambda$, say) to retain some of its
disequilibrium, then the activity ratios can be back-calculated to the
time of isotopic closure (assuming subsequent closed-system
behaviour).  This strategy applies to
$^{234}$U/$^{238}$U-disequilibrium and, for very young samples, to
$^{230}$Th/$^{238}$U-disequilibrium. The complex evolution of the
U-decay series was first described by \citet{bateman1908} and
subsequently applied to U--Pb geochronology by \citet{ludwig1977},
\citet{wendt1985} and \citet{engel2019}. Here we opt for an
alternative formulation, using matrix exponentials
\citep{albarede1995}.  For example, the
\textsuperscript{238}U$\rightarrow$\textsuperscript{206}Pb decay chain
can be expressed in matrix form as follows:
\begin{equation}
  \frac{\partial}{\partial{t}}\left[
    \begin{array}{@{}c@{}}
      n_{38}\\
      n_{34}\\
      n_{30}\\
      n_{26}\\
      n_{06}
    \end{array}
  \right]
  =
  \left[
    \begin{array}{ccccc}
      -\lambda_{38} & 0 & 0 & 0 & 0 \\
      \lambda_{38} & -\lambda_{34} & 0 & 0 & 0 \\
      0 & \lambda_{34} & -\lambda_{30} & 0 & 0 \\
      0 & 0 & \lambda_{30} & -\lambda_{26} & 0 \\
      0 & 0 & 0 & \lambda_{26} & 0\\
    \end{array}
    \right]
  \left[
    \begin{array}{@{}c@{}}
      n_{38}\\
      n_{34}\\
      n_{30}\\
      n_{26}\\
      n_{06}
    \end{array}
    \right]
  \label{eq:matrixdecay}
\end{equation}

\noindent where $\lambda_{26}=0.4332(19)\,\mbox{kyr}^{-1}$
\citep{audi2003}, and the shortest lived intermediate daughters
($<$\,1\,kyr half lives) have been omitted. The solution to
Equation~\ref{eq:matrixdecay} is a so-called matrix exponential:
\begin{equation}
  \left[
    \begin{array}{@{}c@{}}
      n_{38}\\
      n_{34}\\
      n_{30}\\
      n_{26}\\
      n_{06}
    \end{array}
  \right]
  =
  \mbox{expm}\left(
  \left[
    \begin{array}{ccccc}
      -\lambda_{38} & 0 & 0 & 0 & 0 \\
      \lambda_{38} & -\lambda_{34} & 0 & 0 & 0 \\
      0 & \lambda_{34} & -\lambda_{30} & 0 & 0 \\
      0 & 0 & \lambda_{30} & -\lambda_{26} & 0 \\
      0 & 0 & 0 & \lambda_{26} & 0\\
    \end{array}
    \right]
  t
  \right)
  \left[
    \begin{array}{@{}c@{}}
      n_{38}\\
      n_{34}\\
      n_{30}\\
      n_{26}\\
      n_{06}
    \end{array}
    \right]_i
  \label{eq:matrixdecaysolution}
\end{equation}

\noindent which expresses the present day amounts of the different
isotopes as a function of the initial amounts. An interesting result
is obtained by setting $t=\infty$ in
Equation~\ref{eq:matrixdecaysolution} to estimate the activity ratio
under secular equilibrium:
\begin{equation}
  [4/8]_\infty = \frac{\lambda_{234}}{\lambda_{234}-\lambda_{238}} = 1.000055
  \label{eq:secular_equilibrium}
\end{equation}

Note that this activity ratio is \emph{not} exactly equal to
unity. This is because 0.0055\% of ${}^{238}$U is lost during
${}^{234}$U's mean lifetime of $1/\lambda_{34} =
354\,\mbox{kyr}$. Equation~\ref{eq:matrixdecaysolution} can also be
inverted to express the initial amounts as a function of the present
day amounts:
\begin{equation}
  \left[
    \begin{array}{@{}c@{}}
      n_{38}\\
      n_{34}\\
      n_{30}\\
      n_{26}\\
      n_{06}
    \end{array}
  \right]_i
  =
  \mbox{expm}\left(
      -
  \left[
    \begin{array}{ccccc}
      -\lambda_{38} & 0 & 0 & 0 & 0 \\
      \lambda_{38} & -\lambda_{34} & 0 & 0 & 0 \\
      0 & \lambda_{34} & -\lambda_{30} & 0 & 0 \\
      0 & 0 & \lambda_{30} & -\lambda_{26} & 0 \\
      0 & 0 & 0 & \lambda_{26} & 0\\
    \end{array}
    \right]
  t
  \right)
  \left[
    \begin{array}{@{}c@{}}
      n_{38}\\
      n_{34}\\
      n_{30}\\
      n_{26}\\
      n_{06}
    \end{array}
    \right]
  \label{eq:matrixdecaysolutioninverse}
\end{equation}

Equations~\ref{eq:matrixdecaysolution} and
\ref{eq:matrixdecaysolutioninverse} can be used to construct a
concordia diagram in the presence of disequilibrium. If measured
activity ratios are used to infer the initial conditions, then the
concordia line terminates where those inferred activity ratios reach
unrealistic values (e.g., $[4/8]_i=500$ and $[0/8]_i=0$;
Figure~\ref{fig:CorchiaMC}a). Beyond ten or so \textsuperscript{234}U
half lives, it becomes very difficult to estimate $[4/8]_i$ from
$[4/8]_m$, and it is even more difficult to quantify the analytical
uncertainty of the disequilibrium correction. In Section~\ref{sec:MC}
we review the current `Monte Carlo' approach to uncertainty estimation
for initial disequilibrium correction and in Section~\ref{sec:Bayes}
we propose an alternative `Bayesian' approach, which offers
significant advantages for samples that are close to secular
equilibrium.

\section{`Monte Carlo' uncertainty estimation}\label{sec:MC}

Existing data processing software for disequilibrium-corrected
${}^{206}$Pb/${}^{238}$U geo\-chro\-nology, such as \texttt{DQPB}
\citep{pollard2023}, estimate the uncertainty of the disequilibrium
correction by Monte Carlo simulation. Given a linear array of isotopic
data in Tera-Wasserburg space (i.e. $n_{07}/n_{06}$
vs. $n_{38}/n_{06}$), paired with an activity ratio measurement
$[4/8]_m$ with standard error $s[4/8]_m$, this approach works as
follows:

\begin{enumerate}
\item\label{it:2} Draw a random value $[4/8]_t$ from a normal
  distribution with mean $[4/8]_m$ and standard deviation
  $s[4/8]_m$.
\item\label{it:3} Fit a straight line to the U--Pb data and find the
  $[4/8]_i$-value and isochron age ($t$) that are consistent with
  both the U--Pb measurements and $[4/8]_t$. In other words, use
  Equation~\ref{eq:matrixdecaysolutioninverse} to estimate $[4/8]_i$
  from $[4/8]_t$ and repeat this for different values of $t$ until
  the linear fit to the U--Pb data is optimised.
\item Repeat steps~\ref{it:2} and \ref{it:3} until the entire
  distribution of $[4/8]_t$-values has been sampled.
\item If step \ref{it:3} fails or produces physically impossible
  results (e.g., $t<0$), then ignore the corresponding
  $[4/8]_t$-value. Otherwise add the $[4/8]_i$ and $t$-values to a
  list of acceptable results.\label{it:4}
\item Use the spread of the acceptable $[4/8]_i$ and $t$-values to
  quantify their respective uncertainties.
\end{enumerate}

For the purpose of the present study, we have implemented our own
version of this algorithm, using \texttt{R} and \texttt{IsoplotR}
\citep{vermeesch2018c}.  The only major difference between our code
and \texttt{DQPB} is that it does not sample the
$[4/8]_m$-distribution randomly, but uses a targeted approach to
sample $[4/8]_m$ as a sequence of regularly spaced normal
quantiles. This is faster and produces deterministic results that do
not depend on the seed of a random number
generator. Figure~\ref{fig:CorchiaMC} summarises the application of
this approach using the `Corchia' dataset of \citet{pollard2023},
producing identical results to \texttt{DQPB}. To reflect this
equivalence of outcomes, we will refer to our version of the algorithm
as a `Monte Carlo' method, despite the fact that it does not actually
use a random number generator.\medskip

%\begin{figure}[!ht]
{  \centering
  \includegraphics[width=12cm]{figures/Fig1.pdf}
  \captionof{figure}{Output of the `Monte Carlo' algorithm for the
    Corchia dataset of \citet{pollard2023}. a) Tera-Wasserburg
    concordia diagram with disequilibrium-corrected isochron
    ($t=0.5804\,\pm\,0.0086$\,Ma); b) 50 representative samples from the
    distribution of $^{234}$U/$^{238}$U activity ratio measurements;
    c) The corresponding initial $^{234}$U/$^{238}$U activity ratios;
    d) The isochron ages corresponding to the $[4/8]_i$ values
    presented in panel~c.\medskip}
  \label{fig:CorchiaMC}
}%\end{figure}

Next, let us apply the same approach to older materials such as sample
AV03 (Bolt's Farm, South Africa) of \citet{pickering2019}. The
uncorrected U--Pb isochron age for this sample is ${5.6}\pm{0.9}$\,Ma,
which is 22 half-lives of $^{234}$U. Consequently, the measured
present-day $^{234}$U/$^{238}$U activity ratio is statistically
indistinguishable from secular equilibrium, at $[4/8]_m =
{1.0046}\pm{0.0063}$. Despite the lack of measurable disequilibrium,
the `Monte Carlo' approach appears to have successfully applied a
disequilibrium correction, resulting in a corrected age that is less
than half the uncorrected age, with a precision of better than 12\%
(Figure~\ref{fig:AV03MC1}). How is this possible? The answer lies in
the rejected solutions (step \ref{it:4} of the algorithm), which are
marked in black in Figure~\ref{fig:AV03MC1}b. Ignoring these
`physically impossible' initial ratios suppresses the equilibrium
solutions and skews the distribution of Monte Carlo solutions towards
high $[4/8]_i$-values (Figure~\ref{fig:AV03MC1}c) and young ages
(Figure~\ref{fig:AV03MC1}d).\medskip

%\begin{figure}[!ht]
{ \centering
  \includegraphics[width=12cm]{figures/Fig2.pdf}
  \captionof{figure}{Output of the `Monte Carlo' algorithm for sample
    AV03 of \citet{pickering2019}. Panels~a) -- d) are as in
    Figure~\ref{fig:CorchiaMC}. The black dots in panel~b) mark
    synthetic replicates that are rejected because they yield
    physically impossible $[4/8]_i$ and/or $t$-values. The results
    shown in panels~c and d are consistent with the published
    values.\medskip}
  \label{fig:AV03MC1}
}%\end{figure}

To demonstrate that the result of Figure~\ref{fig:AV03MC1} is wrong,
let us replace the measured $^{234}$U/$^{238}$U activity ratio with
the equilibrium ratio (Equation~\ref{eq:secular_equilibrium}):
\[
[4/8]_m = [4/8]_\infty \pm 0.0063
\]

Plugging this value into the `Monte Carlo' algorithm yields an
impossible result (Figure~\ref{fig:AV03MC2}). It has
applied a disequilibrium correction without any actual disequilibrium,
by ignoring exactly half of the $[4/8]_t$ distribution
(Figure~\ref{fig:AV03MC2}b). This was necessary because,
for this old sample, essentially any $[4/8]_t$ value that is less
than the equilibrium ratio would require a negative $[4/8]_i$ ratio,
or a negative isochron age $t$.\medskip

%\begin{figure}[!ht]
{ \centering \includegraphics[width=12cm]{figures/Fig3.pdf}
  \captionof{figure}{The same data as Figure~\ref{fig:AV03MC1}, but
    replacing $[4/8]_m$ with the equilibrium ratio. Note that half
    of the synthetic replicates have been rejected (black
    circles). Even though there is absolutely no evidence for
    disequilibrium, the `Monte Carlo' produces a corrected isochron
    age (panel~d) that is half the uncorrected value (panel~a). This
    result is clearly wrong.\medskip}
  \label{fig:AV03MC2}
}%\end{figure}

\section{A Bayesian approach}\label{sec:Bayes}

The previous section showed that the `Monte Carlo' algorithm produces
incorrect results for samples whose $[4/8]_t$-distributions fall
within, say, three standard errors (i.e., $3\times{s}[4/8]_m$) from
secular equilibrium. One way to address this issue is for the `Monte
Carlo' algorithm to issue a warning when the $[4/8]_m$ value is
close to $[4/8]_\infty$. This is the approach taken by \texttt{DQPB}
\citep{pollard2023}.  In this section we introduce an alternative
approach that automatically handles problematic cases, without the
need to define a nominal `applicability cutoff'. Given any values of
$[4/8]_m$ and $s[4/8]_m$, our new `Bayesian' algorithm proceeds as
follows:

\begin{enumerate}
\item Define a prior distribution for $[4/8]_i$. In a first
  instance, we will use a uniform distribution that stretches between
  a minimum $[4/8]_i$-value $m=0$ and a maximum $[4/8]_i$-value
  $M=20$. However, this uniform distribution can easily be replaced by
  a more informative prior. One flexible way to capture a diversity of
  prior information is the logistic normal distribution:
  \begin{equation}
    \ln\left(\frac{[4/8]_i-m}{M-m}\right) \sim
    \mathcal{N}(\mu,\sigma^2)
    \label{eq:prior}
  \end{equation}
  \noindent where $\mu$ and $\sigma$ are the location and dispersion
  parameters of the distribution, respectively. An application of this
  informative prior will be given in
  Section~\ref{sec:accuracy-precision}.
\item Draw a random sample from the prior distribution, carry out a
  constrained isochron regression and register the resulting age ($t$)
  and corresponding $[4/8]_t$-value. In other words, use
  Equation~\ref{eq:matrixdecaysolution} to estimate $[4/8]_t$ from
  $[4/8]_i$, and repeat this for different values of $t$ until the
  linear fit to the U--Pb data is optimised. Register the likelihood
  of this linear fit using the same algorithm as used for regular
  U--Pb isochron regression \citep{ludwig1998,
    vermeesch2020}.\label{it:2b}
\item Calculate the likelihood of the inferred $[4/8]_t$-values
  under a normal distribution with mean $[4/8]_m$ and standard
  deviation $s[4/8]_m$. Combine with the likelihood of the linear
  fit (obtained in step~\ref{it:2b}) to produce the `posterior'
  probability of initial ratios.\label{it:3b}
\item Repeat steps \ref{it:2b} and \ref{it:3b} to constrain the
  posterior distributions of $[4/8]_i$ and $t$. This can either be
  done using a Markov chain, or with a targeted approach of
  appropriately spaced $[4/8]_i$ values.
\end{enumerate}

Applying this method to the 580\,ka Corchia example
(Figure~\ref{fig:CorchiaBayes}) yields essentially identical results
to the `Monte Carlo' algorithm (Figure~\ref{fig:CorchiaMC}).\medskip

%\begin{figure}[!ht]
{ \centering \includegraphics[width=12cm]{figures/Fig4.pdf}
  \captionof{figure}{The same data as Figure~\ref{fig:CorchiaMC}, but
    using the Bayesian approach. y-axes display the prior probability
    (a), likelihood (b) and posterior probabilities (c and d),
    respectively. For this relatively young sample, the Bayesian
    method yields similar results to the `Monte Carlo'
    solution.\medskip}
  \label{fig:CorchiaBayes}
}%\end{figure}

However, when the Bayesian approach is applied to the older Bolt's
Farm data (Figure~\ref{fig:AV03Bayes}), it produces a very different
result than the `Monte Carlo' approach (Figure~\ref{fig:AV03MC1}). The
posterior distributions for Bolt's Farm sample AV03 (shown in
Figure~\ref{fig:AV03Bayes}c and d) still have maxima at $[4/8]_i=9$
and $t=2.6$\,Ma, just like the `Monte Carlo' distributions
(Figures~\ref{fig:AV03MC1}c and d). But unlike the `Monte Carlo'
solution, the result of the Bayesian approach also assigns a
significant probability to older ages, including the uncorrected U--Pb
date of 5.6\,Ma. The similarity of this posterior distribution to the
prior distribution reflects the fact that the measured
$^{234}$U/$^{238}$U activity ratio contains relatively little
information. The resulting uncertainties are large, but correctly
reflect our ignorance about the true extent of the disequilibrium in
this case.\medskip

Finally, changing the $[4/8]_m$-ratio to the equilibrium value
(Figure~\ref{fig:AV03Bayes}) produces a posterior distribution that is
nearly identical to the prior distribution. This means that the
likelihood function contains almost no information. In other words:
the measured $^{234}$U/$^{238}$U activity ratio does not tell us
anything about the initial disequilibrium (except that
$[4/8]_i<10$).  If it cannot be ruled out that the sample may have
experienced extreme $^{234}$U/$^{238}$U disequilibrium, then it is not
possible to undo the effects of this disequilibrium using the modern
(measured) $^{234}$U/$^{238}$U activity ratio.

%\begin{figure}[!ht]
{ \centering
  \includegraphics[width=12cm]{figures/Fig5.pdf}
  \captionof{figure}{Posterior distributions of $[4/8]_i$ and $t$
    for sample AV03 of \citet{pickering2019}. Panels~a and b represent
    the original U--Pb data of Figure~\ref{fig:AV03MC1}, whereas
    panels~c and d show the modified data of
    Figure~\ref{fig:AV03MC2}. Likelihood functions are provided in
    panels~b of the latter two figures. A uniform prior was used but
    is not shown. The modes of the posterior distributions agree with
    the modes of the `Monte Carlo' solutions. However, whereas the
    `Monte Carlo' algorithm suggests a high degree of confidence in
    the disequilibrium correction, the Bayesian approach shows that
    one cannot rule out a much higher age of the sample, including the
    uncorrected date of 5.6\,Ma.\medskip}
  \label{fig:AV03Bayes}
}%\end{figure}

\section{The case against ${}^{206}$Pb/${}^{238}$U-dating of old carbonates}
\label{sec:case-against}

The applicability range of the $^{206}$Pb/$^{238}$U method depends on
the $[4/8]_i$-ratio and on the precision of the
$[4/8]_m$-mea\-su\-re\-ments.  Although sub-permil level analytical
uncertainties can be routinely achieved for individual
$[4/8]_m$-measurements, the external reproducibility is likely worse
than this in most samples. This is due to a combination of two
competing factors. First, $[4/8]_i$ is negatively correlated with U
concentration \citep{osmond1976,zhou2005,kuribayashi2025}. Second, the
U concentration must exhibit large variations to form a statistically
robust \textsuperscript{206}Pb/\textsuperscript{238}U
isochron.\medskip

The combination of these two effects has the potential to cause
intra-sample variations in $[4/8]_m$ exceeding the analytical
uncertainties. The exact magnitude of the dispersion is unknown
because most speleothem U--Pb dating studies report only one or a few
$[4/8]_m$-values per isochron. Here we will assume a conservative
value of $s[4/8]_m=2\,\permil$, based on a collection of twelve
speleothems analysed by \citet{walker2006}.\medskip

Disequilibrium corrections using measured $^{234}$U/$^{238}$U activity
ratios are only feasible if those activity ratios are statistically
distinguishable from secular equilibrium. To turn these conclusions
into quantitative guidelines, let us define ``statistically
distinguishable'' as ``at least $3\times{s[4/8]_m}$ removed from
secular equilibrium''. Using this definition, a sample with
$[4/8]_i<2.7$ would become indistinguishable from secular
equilibrium after c.\,2\,Ma.  In other words, it is impossible to
correct a 2\,Ma sample whose $[4/8]_i=2.7$, say. The uncorrected
$^{206}$Pb/$^{238}$U isochron age of such a sample would be 2.6\,Ma,
corresponding to a bias of 30\%. Table~\ref{tab:0638sensitivity} shows
the outcomes of the same exercise for a range of other ages.\medskip

%\begin{table}[!ht]
{  \centering
  \captionof{table}{Sensitivity test of the ${}^{206}$Pb/${}^{238}$U method using selected ages.}
  \begin{tabular}{r|cccccc}
    true age (Ma) & 0.5 & 1 & 1.5 & 2 & 2.5 & 3.0 \\
    minimum resolvable $[4/8]_i$& 1.02 & 1.10 & 1.41 & 2.68 & 7.89 & 29.2\\
    maximum bias (\%) & 1.1 & 3.2 & 9.5 & 30 & 97 & 333\\
  \end{tabular}
  \label{tab:0638sensitivity}\medskip
}%\end{table}

Figure~\ref{fig:nomogram} presents a more extensive exploration of the
magnitude (panel~a) and precision (panel~b) of
disequilibrium-corrected $^{206}$Pb/$^{238}$U geochronology assuming
the aforementioned $2\,\permil$ reproducibility. Alternative versions
of this diagram can be generated by modifying the reproducibility
value in the \texttt{R}-code provided in the supplementary
information \citep{vermeesch2025z}.\medskip

Based on these considerations, we judge carbonate $^{206}$Pb/$^{238}$U
geochronology to be unreliable beyond c.\,1.5\,Ma, and impossible
beyond c.\,2\,Ma unless initial $^{234}$U/$^{238}$U disequilibrium can
be confidently ruled out (Figure~\ref{fig:nomogram}).  However, there
is a solution to the conundrum of
$^{234}$U/$^{238}$U-disequilibrium. This solution is the
${}^{207}$Pb/${}^{235}$U isochron method
\citep{richards1998,engel2019,vaks2020,vermeesch2025}.\medskip

\noindent\includegraphics[width=14cm]{figures/Fig6.pdf}
\captionof{figure}{
  Nomogram to assess the applicability of the
  \textsuperscript{206}Pb/\textsuperscript{238}U method in the
  presence of different degrees of initial
  \textsuperscript{234}U/\textsuperscript{238}U-disequilibrium.
  Panels a and b visualise the magnitude and the precision of the
  disequilibrium correction, respectively, for selected values of the
  initial activity ratios $[4/8]_i$.  $t_r$ is the uncorrected date,
  assuming secular equilibrium. $t_c$ is the true age, which equals
  the disequilibrium-corrected date using the expected $[4/8]_m$
  value. $t_u$ and $t_l$ are the disequilibrium-corrected dates using
  present-day \textsuperscript{234}U/\textsuperscript{238}U activity
  ratios of $[4/8]_m + 2\,\permil$ and $[4/8]_m - 2\,\permil$,
  respectively. The dashed line in panel b) marks the relative
  uncertainty interval when no disequilibrium measurement is
  available, defined as the difference between corrected dates 
  assuming initial ${}^{234}$U${/}^{238}$U activity ratios of 1 and 12.}
\label{fig:nomogram}

\section{A potential ${}^{207}$Pb/${}^{235}$U fix to ${}^{206}$Pb/${}^{238}$U's problems}\label{sec:207Pb235U}

In the previous section, we showed that the ${}^{206}$Pb/${}^{238}$U
method's accuracy is hampered by the extreme enrichment (up to double
the equilibrium value or more) of ${}^{234}$U observed in certain
groundwaters
\citep[e.g.,][]{osmond1976,kronfeld1994,kuribayashi2025}. This problem
can be solved by avoiding ${}^{234}$U altogether and sidestepping the
${}^{238}$U--${}^{206}$Pb decay chain in favour of the
${}^{235}$U--${}^{207}$Pb decay chain \citep{neymark2008}.\medskip

There are two kinds of ${}^{207}$Pb/${}^{235}$U-isochrons. The
simplest kind plots ${}^{204}$Pb/${}^{207}$Pb-ratios against
${}^{204}$Pb/${}^{235}$U-ratios, defining the following linear
relationship:
\begin{equation}
  \left[\frac{{}^{204}\mbox{Pb}}{{}^{207}\mbox{Pb}}\right] =
  \left[\frac{{}^{204}\mbox{Pb}}{{}^{207}\mbox{Pb}}\right]_i \left\{ 1
  - \left[\frac{{}^{235}\mbox{U}}{{}^{207}\mbox{Pb}}\right](\exp[\lambda_{35}t]-1) \right\}
  \label{eq:204isochron}
\end{equation}

\noindent where ${}^{204}$Pb is used as a proxy for common
Pb. Alternatively, one can also use ${}^{208}$Pb to fulfil this
role. This gives rise to a ${}^{208}$Pb$_i/{}^{207}$Pb
vs. ${}^{235}$U/${}^{207}$Pb isochron, where ${}^{208}$Pb$_i$ is the
non-radiogenic ${}^{208}$Pb-component (with the decay products of
${}^{232}$Th removed). Equation~\ref{eq:204isochron} is then replaced
with:
\begin{equation}
  \left[\frac{{}^{208}\mbox{Pb}_i}{{}^{207}\mbox{Pb}}\right] =
  \left[\frac{{}^{208}\mbox{Pb}}{{}^{207}\mbox{Pb}}\right]_i \left\{ 1 -
  \left[\frac{{}^{235}\mbox{U}}{{}^{207}\mbox{Pb}}\right](\exp[\lambda_{35}t]-1) \right\}
  \label{eq:208isochron}
\end{equation}

\noindent where
\begin{equation}
  \left[\frac{{}^{208}\mbox{Pb}_i}{{}^{207}\mbox{Pb}}\right] =
  \left[\frac{{}^{208}\mbox{Pb}}{{}^{207}\mbox{Pb}}\right] -
  \left[\frac{{}^{232}\mbox{Th}}{{}^{207}\mbox{Pb}}\right] \exp[\lambda_{32}t] + 1
\end{equation}

\noindent in which $\lambda_{32}=0.0495(25)\,\mbox{Gyr}^{-1}$
\citep{leroux1963} and the ${}^{232}$Th/${}^{207}$Pb-ratio can be
obtained from the product of the ${}^{232}$Th/${}^{238}$U,
${}^{238}$U/${}^{235}$U and ${}^{232}$Th/${}^{207}$Pb ratios. Because
Th is insoluble in water, radiogenic ${}^{208}$Pb is often absent from
carbonates. Therefore, it is generally safe to assume that
$\left[{}^{208}\mbox{Pb}_i/^{207}\mbox{Pb}\right] \approx
\left[{}^{208}\mbox{Pb}/^{207}\mbox{Pb}\right]$.\medskip

The ${}^{235}$U$\rightarrow{}^{207}$Pb method has just one long-lived
intermediate daughter, ${}^{231}$Pa, which requires a correction. Due
to ${}^{231}$Pa's short half-life of 32.65~kyr
\citep[$\lambda_{31}=21.158(71)\,\mbox{Myr}^{-1}$;][]{audi2003}, it is
generally not possible to measure any remaining disequilibrium in the
Myr time range where the ${}^{235}$U$\rightarrow{}^{207}$Pb method
offers a tangible advantage over the
${}^{238}$U$\rightarrow{}^{206}$Pb method. Therefore, the
$[1/5]_i$-value must be assumed.\medskip

Pa is chemically similar to Th, and insoluble in water.  Therefore,
${}^{231}$Pa is always depleted relative to ${}^{235}$U in
carbonates. So whereas $[4/8]_i$ can vary anywhere between 0 and 12
\citep{osmond1976}, $[1/5]_i$ is always less than 1 and can be
safely assumed to be zero. In a worst case scenario, in which one
assumes $[1/5]_i=0$ but the true activity ratio is $[1/5]_i=1$,
this would only bias the ${}^{235}$Pb/${}^{235}$U age by a relatively
small amount (Table~\ref{tab:0735sensitivity}).\medskip

%\begin{table}[!ht]
{ \centering \captionof{table}{Sensitivity test of the
    ${}^{207}$Pb/${}^{235}$U method against $^{231}$Pa
    disequilibrium.}
  \begin{tabular}{r|cccccc}
    true age (Ma) & 0.5 & 1 & 1.5 & 2 & 2.5 & 3.0 \\
    maximum bias (\%) & 9.4 & 4.7 & 3.1 & 2.4 & 1.9 & 1.6
  \end{tabular}
  \label{tab:0735sensitivity}\medskip
}%\end{table}

The degree of potential bias of the ${}^{207}$Pb/${}^{235}$U method
decreases with increasing age, unlike the ${}^{206}$Pb/${}^{238}$U
method, whose bias increases with age
(Table~\ref{tab:0638sensitivity}). In this sense, the
${}^{207}$Pb/${}^{235}$U and ${}^{206}$Pb/${}^{238}$U methods are
complementary to each other. The ${}^{207}$Pb/${}^{235}$U is most
accurate for samples older than 1\,Ma, whereas the
${}^{206}$Pb/${}^{238}$U is more accurate for samples younger than
1\,Ma. Note that the latter is similar to the applicability range of
the ${}^{230}$Th/U method. So one could argue that
disequilibrium-corrected ${}^{206}$Pb/${}^{238}$U dating is of limited
use to carbonate U--Pb geochronology \citep[except to infer
  {[4/8]\textsubscript{\it{i}}};][]{engel2019}. Although the
${}^{207}$Pb/${}^{235}$U method outperforms the
${}^{206}$Pb/${}^{238}$U method at c.\,1\,Ma in terms of accuracy, its
poorer precision means that its potential benefits do not materialise
until c.\,2\,Ma. In the next Section, we will demonstrate this by
applying both methods to three different case studies.

\section{Case studies}
\label{sec:accuracy-precision}

Having made a largely theoretical case against
\textsuperscript{206}Pb/\textsuperscript{238}U dating and for
\textsuperscript{207}Pb/\textsuperscript{235}U dating of old
carbonates that are suspected to have experienced initial
\textsuperscript{234}U/\textsuperscript{238}U disequilibrium, we will
now compare and contrast the two chronometers using three practical
case studies.  The first example will demonstrate the accuracy of the
\textsuperscript{207}Pb/\textsuperscript{235}U method by showing its
consistency with disequilibrium-corrected
\textsuperscript{206}Pb/\textsuperscript{238}U dates of young
($<2$\,Ma) samples.\medskip

The second example uses ID-TIMS data to serve two purposes. First, it
will show that the \textsuperscript{207}Pb/\textsuperscript{235}U
method produces more accurate, more consistent, and more precise
results than the \textsuperscript{206}Pb/\textsuperscript{238}U method
for older ($>2$\,Ma) carbonates. Second, it will demonstrate how the
Bayesian framework can use prior information to overcome the the
inaccuracy of the \textsuperscript{206}Pb/\textsuperscript{238}U
method.\medskip

In the third case study we apply the
\textsuperscript{207}Pb/\textsuperscript{235}U method to LA-ICP-MS
data using ${}^{208}$Pb as a proxy for common Pb.  In addition to
highlighting a successful application where the
\textsuperscript{207}Pb/\textsuperscript{235}U method produces
demonstrably superior results to the
\textsuperscript{206}Pb/\textsuperscript{238}U method, this dataset
also illustrates limitations of the
\textsuperscript{207}Pb/\textsuperscript{235}U approach with a sample
that yields a precise \textsuperscript{206}Pb/\textsuperscript{238}U
isochron and an unusably imprecise
\textsuperscript{207}Pb/\textsuperscript{235}U isochron.

\subsection{ID-ICP-MS data from Siberia}
\label{sec:Vaks}

A rich dataset of 72 speleothem dates is available from the Botovskaya
and Ledyanaya Lenskaya (LLC) caves in Siberia. \citet{vaks2020} used
the U--Pb method to extend an important palaeoclimatological archive
from these caves that was previously dated using the ${}^{230}$Th--U
disequilibrium method \citep{vaks2013}.  Samples were analysed by
isotope dilution ICP-MS and were found to exhibit a significant level
of ${}^{234}$U/${}^{238}$U disequilibrium.  U--Pb ages were estimated
using a two-step procedure.  First, the common-Pb contribution was
removed by two-point isochron regression through an inherited
composition that was inferred by inspection of apparent linear trends
in Tera-Wasserburg concordia space. Second, a ${}^{234}$U/${}^{238}$U
disequilibrium correction was applied to the radiogenic end-member
composition, using the procedures described in
Section~\ref{sec:diseq}.  This correction combined the measured
${}^{234}$U/${}^{238}$U activity ratios with an assumed absence of
initial ${}^{230}$Th and ${}^{231}$Pa.\medskip

The inferred $[4/8]_i$-values were c.\,2 and 3--5 for LLC and
Botovskaya cave, respectively. This corresponds to an age correction
of 15\% for LLC and 60\% for Botovskaya cave
(Figure~\ref{fig:Siberia}). Uncertainties were estimated using the
Bayesian procedure of Section~\ref{sec:Bayes}, making the optimistic
assumption that the analytical uncertainty of the
$[4/8]_m$-measurements faithfully captures all sources of
dispersion. The scatter of the $[4/8]_i$-values suggests that this
may not be the case. This caveat notwithstanding, the
disequilibrium-corrected ${}^{206}$Pb/${}^{238}$U and
${}^{207}$Pb/${}^{235}$U ages overlap within uncertainty in all but
four of the samples.\medskip

The ${}^{207}$Pb/${}^{235}$U age uncertainty is invariably larger than
the ${}^{206}$Pb/${}^{238}$U age uncertainty.  In fact, below
c.\,1\,Ma, it could be argued that the ${}^{207}$Pb/${}^{235}$U age
uncertainties are unusably imprecise ($s[t]/t>50\%$).  However, above
c.\,1\,Ma, the uncertainty reduces to acceptable levels
($s[t]/t<5\%$).  Extrapolating this trend further into the past
confirms the earlier assertion that beyond c.\,2\,Ma, the
${}^{207}$Pb/${}^{235}$U method outperforms the
${}^{206}$Pb/${}^{238}$U method in both accuracy and precision.

{ \centering \includegraphics[width=12cm]{figures/Fig7.pdf}
  \captionof{figure}{ Reanalysis of the speleothem data of
    \citet{vaks2020}. This record stacks together several speleothems,
    which are arranged in stratigraphic order for each cave. a)
    uncorrected (circles) and corrected (black error bars)
    ${}^{206}$Pb/${}^{238}$U dates, juxtaposed next to the
    ${}^{207}$Pb/${}^{235}$U dates (blue error bars) for the same
    samples. b) measured ($[4/8]_m$, circles) and inferred initial
    ($[4/8]_i$, error bars) ${}^{234}$U/${}^{238}$U-activity
    ratios. All error bars represent Bayesian 95\% credible
    intervals. Sample~5 does not have a $[4/8]_m$ measurement and
    was assumed to be in secular equilibrium. Sample~41 is an outlier
    that has an anomalously high common Pb concentration and is only
    included for the sake of completeness.\medskip}
  \label{fig:Siberia}
}

\subsection{ID-TIMS data for ASH-15}
\label{sec:ASH15}

ASH-15 is a carbonate U--Pb dating reference material sourced from a
flowstone in Ashalim cave of southern Israel \citep{nuriel2021}.
37~ID-TIMS measurements were obtained from the flowstone, including 12
from horizon D and 25 from horizon K. The latter are shown in
Figure~\ref{fig:ASH15}. \citet{nuriel2021} report an uncorrected
semitotal-Pb/U isochron age of $2.965 \pm 0.011$\,Ma for
ASH-15.\medskip

No \textsuperscript{234}U/\textsuperscript{238}U activity ratio
measurements are available for ASH-15D and ASH-15K. However, two other
horizons of the same flowstone (ASH-15A+B and ASH-15C1) are
characterised by $[4/8]_m$-values of ${0.99939}\pm{0.00108}$ and
${0.99925}\pm{0.0015}$, indistinguishable from secular
equilibrium. Younger flowstones in Ashalim cave yield an average
$[4/8]_i$-value of 1.0470 with a standard deviation of 0.01492
\citep{vaks2010,vaks2013b}. A wider survey of 904 speleothem samples
dated in southern and central Israel by \citet{chaldekas2022} have
average $[4/8]_i$-values of $1.081\pm{0.138}$. Despite this lack of
observable disequilibrium, \citet{mason2013} suggest a
$[4/8]_i$-value of 1.5--2.0 to explain the minor degree of
discordance of the common-Pb corrected Tera-Wasserburg ratios.\medskip

To investigate the effect of initial U-series disequilibrium on
ASH-15, Figures~\ref{fig:ASH15}a and b apply the Bayesian inversion
algorithm to the ASH-15K data, using the $[4/8]_m$-value of ASH-15C1
and assuming that $[0/8]_i=0$ (i.e. no initial
\textsuperscript{230}Th). In a first attempt, we will use the same
uniform prior from $m=0$ to $M=20$ as before. This results in a
disequilibrium-corrected
\textsuperscript{206}Pb/\,\textsuperscript{238}U-isochron age of
$3.47+0.053/-0.794$\,Ma, and an inferred $[4/8]_i$ ratio of
$0.067+2.23/-0.052$. As expected, initial equilibrium is very
plausible with $P([4/8]_i>1.0)=0.47$. The initial activity ratio
preferred by \citet{mason2013} cannot be ruled out either, but is less
likely, with $P([4/8]_i>1.5)=0.25$ and $P([4/8]_i>2.0)=0.09$.\medskip

In a second attempt, we used the $[4/8]_i$-values of
\citet{chaldekas2022} to construct an informative prior, using the
logistic normal formulation of Equation~\ref{eq:prior} with $m=1$,
$M=3$, $\mu=1.081$ and $\sigma=0.2$.  This produces a posterior
distribution for $[4/8]_i$ that is, essentially, identical to the
prior, confirming that the $[4/8]_m$ data carry virtually no
additional information. The corresponding disequilibrium-corrected
\textsuperscript{206}Pb/\,\textsuperscript{238}U-isochron age is
$3.107\pm{0.065}$\,Ma.\medskip

In contrast with the widely varying scenarios for the
\textsuperscript{206}Pb/\,\textsuperscript{238}U method, the
\textsuperscript{207}Pb/\,\textsuperscript{235}U isochron age
calculation is straightforward:

\begin{enumerate}
\item An uncorrected \textsuperscript{207}Pb/\,\textsuperscript{235}U
  isochron age of $3.039\pm0.068$\,Ma, assuming $[1/5]_i\approx{1}$.
\item A corrected \textsuperscript{207}Pb/\,\textsuperscript{235}U
  isochron age of $3.086\pm0.068$\,Ma, assuming $[1/5]_i=0$.
\end{enumerate}

\noindent which are nearly identical to the
\textsuperscript{206}Pb/\textsuperscript{238}U date using the
informative $[4/8]_i$-prior. We would like to conclude the
discussion of ASH-15 by remarking that disequilibrium issues do not
affect the suitability of this sample as a reference material for
in-situ U--Pb geochronology. This is because standardisation is done
relative to the \emph{uncorrected} isotopic composition
\citep{horstwood2016}.\medskip

{ \centering 
  \includegraphics[width=\linewidth]{figures/Fig8.pdf}
  \captionof{figure}{
    ID-TIMS U--Pb data for flowstone ASH-15K \citep{nuriel2021}. a)
    informative prior for $[4/8]_i$, based on the data compilation of
    \citet{chaldekas2022}; b) likelihood function for $[4/8]_m$, using
    the activity ratio measurements of ASH-15C1 \citep{vaks2013}; c)
    and d) posterior probabilities for $[4/8]_i$ and $t$,
    respectively; e) the ${}^{206}$Pb/${}^{238}$U isochron, fitted
    using the model-3 algorithm of \citet{vermeesch2024}, with the
    excess dispersion shown as a horizontal 95\% error bar; and f) the
    model-1 ${}^{207}$Pb/${}^{235}$U isochron. The grey uncertainty
    bands represent the standard errors of the isochron fits and do
    not reflect the Bayesian credible intervals.}
  \label{fig:ASH15}
}

\subsection{LA-ICP-MS data from Siberia}

For the final case study, we return from Israel to the Botovskaya cave
deposits in Siberia. Section~\ref{sec:Vaks} and
Figure~\ref{fig:Siberia} show abundant evidence that this cave is
strongly enriched in initial ${}^{234}$U, with $[4/8]_i$-values
ranging from 3 to 5 according to the results of \citet{vaks2020}. The
effect of this strong disequilibrium can be confidently undone for the
young ($<500$\,ka) speleothems shown in
Figure~\ref{fig:Siberia}. However, Botovskaya cave also contains
speleothems that are considerably older than this, going all the way
back to the Plio-Pleistocene. By now it should be clear that the
${}^{206}$Pb/${}^{238}$U method is ill suited to unlock this older
archive. Figure~\ref{fig:Botovskaya} summarises some preliminary U--Pb
results from two of these older cave deposits (sample SB-1625-22 and
sample SB-72-8) obtained by LA-ICP-MS.\medskip

As mentioned before under the discussion of the Bolt's Cave data, the
${}^{204}$Pb measurements produced by this technique are imprecise and
potentially inaccurate. In this case, ${}^{208}$Pb was measured and
can be used as a substitute for ${}^{204}$Pb.  Th/U ratios (also
measured by LA-ICP-MS) were extremely low, allowing us to ignore the
radiogenic ${}^{208}$Pb contribution.\medskip

The uncorrected ${}^{206}$Pb/${}^{238}$U isochron age of sample
SB-1625-22 is $2.66\pm{0.10}$\,Ma and exhibits significant
overdispersion with respect to the analytical uncertainties (MSWD=36).
Model-3 isochron regression \citep[\emph{sensu}][]{vermeesch2024}
indicates that this excess scatter is equivalent to an age dispersion
of $104\pm{28}$\,kyr (Figure~\ref{fig:Botovskaya}a). In reality, the
excess dispersion around the isochron is unlikely to reflect
diachronous isotopic closure. A more likely explanation for the
scatter of the ${}^{206}$Pb/${}^{238}$U data is spatial variability of
the $[4/8]_i$-values, as discussed in
Section~\ref{sec:case-against}. In summary, the actual dispersion
estimate probably has no physical meaning, but the isochron age should
be as accurate as mathematically possible. The reduction of the
scatter around the Botovskaya isochrons (Figures~\ref{fig:Botovskaya}a
and c) towards the y-intercept also suggests that the common Pb-ratio
is not substantially correlated with the postulated heterogeneous
initial disequilibrium.\medskip

Given the antiquity of the sample and the difficulty of measuring
$[4/8]_m$ by LA-ICP-MS, no initial disequilibrium measurement was
made. Switching from ${}^{206}$U/${}^{238}$U to
${}^{207}$Pb/${}^{235}$U isochron space lowers the age to
$1.60\pm{0.10}$\,Ma whilst reducing the dispersion of the data around
the isochron line (MSWD=1.5, Figure~\ref{fig:Botovskaya}b). To verify
the accuracy of this result, it is useful to point out that a
$[4/8]_i$-value of 3.9 would bring the corrected
${}^{206}$Pb/${}^{238}$U isochron in alignment with the
${}^{207}$Pb/${}^{235}$U isochron. Such a value is consistent with the
initial activity ratios of the more recent Botovskaya deposits
(Figure~\ref{fig:Siberia}b). This not only supports the accuracy of
the ${}^{207}$Pb/${}^{235}$U isochron results, but also suggests that
the $[4/8]_i$ ratios have remained stable over hundreds of thousands
of years.\medskip

We would like to conclude the results section by drawing attention to
the fact that the ${}^{207}$Pb/${}^{235}$U method is not always
successful. Figures~\ref{fig:Botovskaya}c and d show that Botovskaya
sample SB-72-8 produces a well defined linear array in
${}^{206}$Pb/${}^{238}$U isochron space, but fails to do so in
${}^{207}$Pb/${}^{235}$U isochron space.  Such cases are not rare. The
${}^{207}$Pb/${}^{235}$U approach only works in samples that are
sufficiently rich in U and sufficiently poor in common Pb. Speleothems
from Botovskaya are rich in U (30--170\,ppm for SB-72-8), so that here
the problem seems to originate from common Pb (0.4--4\,ppm for
SB-72-8).\medskip

{ \centering
  \includegraphics[width=.7\linewidth]{figures/Fig9.pdf}
  \captionof{figure}{LA-ICP-MS data for two speleothems from 
    Botovskaya cave in Siberia. a) model-3 ${}^{206}$Pb/${}^{238}$U
    isochron regression for sample SB-1625-22. The equivalent model-1
    isochron age (with MSWD=9) is $2.65\pm{0.10}$\,Ma. b) model-1
    ${}^{207}$Pb/${}^{235}$U isochron for SB-1625-22; c) model-1
    ${}^{206}$Pb/${}^{238}$U and d) ${}^{207}$Pb/${}^{235}$U isochrons
    for sample SB-72-8 ($2.73\pm{1.33}$\,Ma).}
  \label{fig:Botovskaya}
}

\section{Implementation in \texttt{IsoplotR}}\label{sec:IsoplotR}

All the methods described in this paper have been implemented in the
\texttt{IsoplotR} toolbox for geochronological data processing
\citep{vermeesch2018c}.  The matrix exponential disequilibrium
correction method of Section~\ref{sec:diseq} has been part of
\texttt{IsoplotR} since version 3.0, whereas the deterministic
Bayesian uncertainty estimation routine of Section~\ref{sec:Bayes} was
introduced in version 5.2.  At the time of writing, \texttt{IsoplotR}
(version 6.7) supports twelve different U--Pb data
formats. Disequilibrium corrected U--Pb isochron regression is
available for all these formats, in different forms.\medskip

Formats~1--3 contain neither ${}^{204}$Pb nor ${}^{208}$Pb. Therefore,
isochron regression for these formats must be done by semitotal-Pb/U
regression in Tera-Wasserburg concordia space. Formats~4--6 include
${}^{204}$Pb as a common Pb tracer. These formats permit the
calculation of both ${}^{206}$Pb/${}^{238}$U and
${}^{207}$Pb/${}^{235}$U isochrons, either jointly \citep[by
  three-dimensional total-Pb/U isochron regression;][]{ludwig1998} or
separately.  To take full advantage of the ${}^{207}$Pb/${}^{235}$U
method's superior accuracy, it is recommended to use the
two-dimensional option. Formats~7 and 8 use ${}^{208}$Pb as a common
Pb tracer. They are also amenable to both ${}^{206}$Pb/${}^{238}$U and
${}^{207}$Pb/${}^{235}$U isochron regression, either jointly \citep[by
  total Pb/U--Th regression;][]{vermeesch2020} or separately.
Formats~9--10 and 11--12 are simplified versions of formats~4--6 and
7--8, respectively, which only permit two-dimensional regression.
Formats~9 and 11 are meant for ${}^{206}$Pb/${}^{238}$U isochron
regression, whereas formats~10 and 12 are meant for
${}^{207}$Pb/${}^{235}$U isochron regression.\medskip

The disequilibrium corrections can be accessed from
\texttt{IsoplotR}'s GUI (either online or offline) by using the
`isochron' function and ticking the `apply disequilibrium correction'
check box in the options menu. Alternatively, the same functionality
can also be accessed from the command-line API. Bayesian uncertainty
estimation is possible using either interface, but visualising the
posterior distributions of the parameter space is currently only
possible from the command line.\medskip

The online supplement provides all the \texttt{R} code that was used
to reproduce the figures in this paper \citep{vermeesch2025z}.

\section{Conclusions}\label{sec:conclusions}

In this paper, we presented a critical appraisal of carbonate U--Pb
geochronology, and proposed three improvements to the
technique. First, we introduced a matrix exponential solution to the
initial disequilibrium problem, extending the work of
\citet{albarede1995}. This formulation produces identical results to
the conventional solution by \citet{engel2019}, but can be written out
more succinctly, and can more easily be modified to suit other
problems. For example, the matrix exponential approach can be adjusted
to calculate disequilibrium-corrected U--Th--He ages
\citep{farley2002b,danisik2017b}.  Second, we presented a
deterministic Bayesian algorithm to quantify the statistical
uncertainty associated with the disequilibrium correction. This
algorithm was used to demonstrate that, for samples older than
c.\,2\,Ma, disequilibrium-corrected ${}^{206}$Pb/${}^{238}$U
geochronology is unreliable. Third, we advocated the use of the
${}^{207}$Pb/${}^{235}$U isochron method as a more accurate
alternative to the ${}^{206}$Pb/${}^{238}$U method.\medskip

Although our findings are most relevant to young carbonates, the
inaccuracy of the ${}^{206}$Pb/${}^{238}$U method equally applies to
old samples.  Only the relative difference between the
${}^{206}$Pb/${}^{238}$U and ${}^{207}$Pb/${}^{235}$U ages reduces
with time. The absolute difference remains constant at up to 4\,Myr
(Equation~\ref{eq:worst-case}). The corresponding systematic
uncertainty cannot be removed without making unverifiable assumptions
about the initial ${}^{234}$U/${}^{238}$U activity ratio.  For samples
older than $>100$\,Ma, say, the systematic error caused by initial
disequilibrium is generally smaller than the random errors associated
with the isotope ratio measurements. However, given a sufficiently
precise set of isochrons, it is theoretically possible to reconstruct
the U-disequilibrium conditions at the time of isotopic closure from
the difference between the ${}^{206}$Pb/${}^{238}$U and
${}^{207}$Pb/${}^{235}$U clocks.\medskip

\citet{engel2019} advocate using the same procedure in Quaternary
studies.  They propose a two-step procedure, whereby the difference
between the ${}^{206}$Pb/${}^{238}$U and ${}^{207}$Pb/${}^{235}$U
dates is used to estimate $[4/8]_i$; and this $[4/8]_i$-value is
then used to calculate a corrected
${}^{206}$Pb/${}^{238}$U-age. \texttt{IsoplotR} implements a one-step
algorithm that achieves the same goal using the total-Pb/U algorithm
of \citet{ludwig1998} and the total-Pb/U--Th algorithm of
\citet{vermeesch2020}. However, we would like to add a note of caution
about the usefulness of this joint regression procedure. Beyond
c.\,2\,Ma, all the age-resolving power of the paired
${}^{206}$Pb/${}^{238}$U and ${}^{207}$Pb/${}^{235}$U approach resides
in the ${}^{207}$Pb/${}^{235}$U clock, so the ${}^{206}$Pb/${}^{238}$U
data add no value.

\subsection*{Code availability}

\texttt{IsoplotR} is free software released under the GPL-3
license. The package and its source code are available from
https://cran.r-project.org/package=IsoplotR \citep[last access:
  July~21, 2025,][]{vermeesch2018c}. The test data used in
Section~\ref{sec:IsoplotR} are provided in the supplementary
information \citep{vermeesch2025z}

\subsection*{Author Contributions}

NM developed the matrix exponential solution to the initial
disequilibrium problem \citep{mclean2016b}; PV developed the
deterministic Bayesian inversion algorithm; RP proposed the
${}^{207}$Pb/${}^{235}$U fix to the ${}^{206}$Pb/${}^{238}$U problems;
AV, TG and SB provided the Siberian data; PV wrote the paper, with
feedback from the other authors.

\subsection*{Competing Interests}

PV and NM are Associate Editors of \emph{Geochronology}.

\subsection*{Acknowledgements}
This research has been supported by the Natural Environment Research
Council (grant no. NE/T001518/1, awarded to PV) and the Leverhulme
Trust (grant no. RPG-20202-334, awarded to SFMB). This paper benefited
from thorough reviews by Perach Nuriel, Timothy Pollard and Robyn
Pickering.

\bibliographystyle{copernicus}
\bibliography{/home/pvermees/Dropbox/biblio.bib}

\end{document}
